\documentclass[11pt]{article}
\usepackage[margin=1in]{geometry}
\usepackage[T1]{fontenc}
\usepackage{lmodern}
\usepackage{microtype}
\usepackage{amsmath}
\usepackage{hyperref}
\ExplSyntaxOn
\ior_new:N \g_binaml_metadata_ior
\tl_new:N \g_binaml_title_tl
\seq_new:N \g_binaml_authors_seq
\seq_new:N \g_binaml_author_emails_seq
\tl_new:N \g_binaml_date_tl
\tl_new:N \g_binaml_status_tl
\cs_new_protected:Npn \binaml_read_metadata:
  {
    \ior_open:Nn \g_binaml_metadata_ior { metadata.yaml }
    \ior_str_map_inline:Nn \g_binaml_metadata_ior
      {
        \regex_match:nnTF { \A title: } { ##1 }
          {
            \tl_gset:Nn \g_binaml_title_tl { ##1 }
            \regex_replace_once:nnN
              { \A title: \s* " } { } \g_binaml_title_tl
            \regex_replace_once:nnN
              { " \s* \Z } { } \g_binaml_title_tl
          }
          {
            \regex_match:nnTF { \A \s* email: } { ##1 }
              {
                \tl_set:Nn \l_tmpa_tl { ##1 }
                \regex_replace_once:nnN
                  { \A \s* email: \s* " } { } \l_tmpa_tl
                \regex_replace_once:nnN
                  { " \s* \Z } { } \l_tmpa_tl
                \seq_gput_right:NV \g_binaml_author_emails_seq \l_tmpa_tl
              }
              {
                \regex_match:nnTF { \A date: } { ##1 }
                  {
                    \tl_gset:Nn \g_binaml_date_tl { ##1 }
                    \regex_replace_once:nnN
                      { \A date: \s* } { } \g_binaml_date_tl
                  }
                  {
                    \regex_match:nnTF { \A status: } { ##1 }
                      {
                        \tl_gset:Nn \g_binaml_status_tl { ##1 }
                        \regex_replace_once:nnN
                          { \A status: \s* } { } \g_binaml_status_tl
                      }
                      {
                        \regex_match:nnT { \A \s* - \s* name: } { ##1 }
                          {
                            \tl_set:Nn \l_tmpa_tl { ##1 }
                            \regex_replace_once:nnN
                              { \A \s* - \s* name: \s* } { } \l_tmpa_tl
                            \seq_gput_right:NV \g_binaml_authors_seq \l_tmpa_tl
                          }
                      }
                  }
              }
          }
      }
    \ior_close:N \g_binaml_metadata_ior
  }
\binaml_read_metadata:
\NewDocumentCommand \BinamlTitle { } { \tl_use:N \g_binaml_title_tl }
\NewDocumentCommand \BinamlAuthors { }
  {
    \seq_map_indexed_inline:Nn \g_binaml_authors_seq
      {
        ##2 \\
        \href
          {mailto:\seq_item:Nn \g_binaml_author_emails_seq { ##1 }}
          {\seq_item:Nn \g_binaml_author_emails_seq { ##1 }}
        \int_compare:nNnT { ##1 } < { \seq_count:N \g_binaml_authors_seq }
          { \and }
      }
  }
\NewDocumentCommand \BinamlDate { }
  {
    \tl_use:N \g_binaml_date_tl
    \str_if_eq:eeT { \tl_use:N \g_binaml_status_tl } { draft }
      { \space (Draft) }
  }
\ExplSyntaxOff
\title{\BinamlTitle}
\author{\BinamlAuthors}
\date{\BinamlDate}
\begin{document}
\maketitle
\begin{abstract}
This paper specifies reproducible synthetic online regression and multiclass
classification tasks for binary feature vectors.
Concept drift shows up across real-world ML, yet most frameworks still assume
stationary train/serve data; these tasks take nonstationarity as a first-class
requirement with controllable drift axes rather than an afterthought.
A synthetic setting makes it easy to experiment with controlled, reproducible
scenarios---known latent mechanisms, seeded trajectories, and explicit scenario
knobs---so adaptation can be studied without confounding from opaque real
streams.  The paper defines both synthetic environments, their independent
sources of concept drift, the model boundary, and the prequential evaluation
protocol for each task.  It is a task and library specification, not a
performance report.
\end{abstract}

\section{Notation}
\begin{description}
  \item[$\xi,d,m,\sigma$] Root seed, input width, number of components, and noise standard deviation.
  \item[$x_t,g_t$] Current binary input and latent activity-gate states.
  \item[$K,y_t$] Number of classes and observed class label (classification).
  \item[$s_{t,c},\tilde s_{t,c}$] Noiseless and noisy class score
  (classification).
  \item[$y_t,\tilde y_t$] Noiseless and observed regression targets.
  \item[$b_t,\epsilon_t,\sigma$] Intercept, observation noise, and its standard deviation.
  \item[$f_i:\{0,1\}^d\to\{0,1\},w_i$] Fixed binary indicator function and its fixed weight.
  \item[$k_i,I_i,\nu_i$] Clause degree, ordered selected-feature indices, and per-literal negation flags.
  \item[$k_{\min},k_{\max},p_{\neg}$] Clause degree bounds and negated-literal probability.
  \item[$q_{\max},p_{\min},p_{\max}$] Maximum parent count and CPT range.
  \item[$p_x,p_{\mathrm{sample,min}}^x,p_{\mathrm{sample,max}}^x$] Input distribution-sampling probability and node-sampling-probability range.
  \item[$p_g,p_{\mathrm{sample,min}}^g,p_{\mathrm{sample,max}}^g$] Gate distribution-sampling probability and node-sampling-probability range.
  \item[$w_{\min},w_{\max},b_{\min},b_{\max}$] Weight and intercept ranges.
  \item[$p_b$] Intercept-sampling probability.
  \item[$S_t^x,S_t^g$] Input- and gate-distribution sampling indicators.
  \item[$S_{t,j}^x,S_{t,j}^g$] Input-node and gate-node sampling indicators.
  \item[$S_t^b$] Intercept-sampling indicator.
\end{description}

\section{Regression task}
At time $t$, an environment emits a binary vector $x_t\in\{0,1\}^d$ and an
observed scalar target $\tilde y_t$.  The environment also maintains binary
activity gates $g_t\in\{0,1\}^m$, one per target component, which are never
supplied to a model.  Keeping $g_i$ hidden keeps activity latent: if gates were
observed, $g_i f_i(x)$ would collapse to just another visible component
function, and the same $f_i$ could not turn on and off without becoming a
distinct observed feature.  A latent target is
\[
y_t=b_t+\sum_{i=1}^{m}w_i g_{t,i} f_i(x_t),
\qquad
\epsilon_t\sim\mathcal N(0,\sigma^2),
\qquad
\tilde y_t=y_t+\epsilon_t,
\]
where $\epsilon_t$ is zero-mean Gaussian noise.  This form asks models to
discover structure in input bits.  Activity gates provide controllable levels
of concept drift without rewriting the underlying components.  Each
$f_i:\{0,1\}^d\to\{0,1\}$ is a binary indicator (boolean component) defined
as one conjunction of literals over a selected subset of input bits: each
literal is either a coordinate or its negation, and all selected literals must
be satisfied for $f_i(x)=1$.

\paragraph{Fixed target construction.}
For each component, draw the clause degree
\[
 k_i\sim\operatorname{DiscreteUniform}
 \{k_{\min},\ldots,k_{\max}\},
\]
sample distinct input indices
\[
 I_i\sim\operatorname{Uniform}\!\left(
 \binom{\{1,\ldots,d\}}{k_i}\right),
\]
and, independently for each selected index, draw a negation flag
\[
 \nu_{i,\ell}\sim\operatorname{Bernoulli}(p_{\neg})
 \qquad (\ell=1,\ldots,k_i).
\]
The component is
\[
 f_i(x)=\bigwedge_{\ell=1}^{k_i}\ell_{i,\ell}(x),
\]
where $\ell_{i,\ell}(x)=x_{I_{i,\ell}}$ if $\nu_{i,\ell}=0$ and
$\ell_{i,\ell}(x)=\neg x_{I_{i,\ell}}$ otherwise.
Weights are sampled once, for example
$w_i\sim\operatorname{Uniform}(w_{\min},w_{\max})$, and remain fixed for the
trajectory.  The initial intercept is
\[
 b_0\sim\operatorname{Uniform}(b_{\min},b_{\max}).
\]
Thus functions and weights are fixed after initialization.  Holding components
fixed keeps drift attributable to gates, intercept, noise, and input-state
dynamics, and lets evaluation ask how fast a model re-adapts to scenarios it
has already faced---for example because it built and kept reusable
representations rather than forgetting them.  Output changes arise from
activity-gate transitions and independent intercept sampling.

\section{Input and gate state dynamics}
Let $X_t=(X_{t,1},\ldots,X_{t,d})$ be the input state and
$G_t=(G_{t,1},\ldots,G_{t,m})$ the activity-gate state.  We describe input
sampling below.  Gate sampling uses the same procedure independently.
\[
 X_t\sim P_x.
\]
The input DAG has $d$ nodes.  At initialization, sample its topological order:
\[
 \pi_x\sim\operatorname{Uniform}(\mathcal S_d).
\]
For each node, uniformly sample a parent subset from all subsets of earlier
nodes in its topological order with at most $q_{\max}$ parents.  Independently
sample every conditional Bernoulli probability in its CPT uniformly from
$[p_{\min},p_{\max}]$.  Sample $X_0$ ancestrally from this DAG, then sample
its shared node-sampling probability:
\[
p_{\mathrm{sample}}^x\sim\operatorname{Uniform}
\left(p_{\mathrm{sample,min}}^x,
      p_{\mathrm{sample,max}}^x\right).
\]
Apply this construction independently to $G_t$: use $m$ nodes,
$\pi_g\sim\operatorname{Uniform}(\mathcal S_m)$, and the gate-specific
$p_g$, $p_{\mathrm{sample,min}}^g$, and $p_{\mathrm{sample,max}}^g$.
Neither topological order alters the model-visible input-coordinate order.

\paragraph{State updates.}
After an emitted sample, draw the input-distribution sampling indicator:
\[
 S_t^x\sim\operatorname{Bernoulli}(p_x).
\]
If $S_t^x=1$, sample a new input DAG order, parents, CPTs,
and $p_{\mathrm{sample}}^x$ while retaining $X_t$.  Then sample every input
node through
\[
 S_{t,j}^x\sim\operatorname{Bernoulli}(p_{\mathrm{sample}}^x)
 \quad (j\in\{1,\ldots,d\}),
\]
then, in topological order, ancestrally sample nodes with indicator one under
the current input DAG.

Each sampled node is drawn from its CPT conditional on its current parent
values.  This partial ancestral sampling intentionally does not guarantee that
the current state follows the configured DAG distribution.  The gate state
uses this same update independently: draw $S_t^g\sim\operatorname{Bernoulli}(p_g)$,
retain $G_t$ if it is zero or sample a new gate distribution if it is one,
then sample gate nodes in order with shared probability
$p_{\mathrm{sample}}^g$.  Each shared node-sampling probability controls
persistence and is sampled when its distribution is initialized or sampled
again.

\section{Stream-generation procedure}
\noindent\textbf{Configuration}
\begin{description}
  \item[\textbf{Run}] Root seed $\xi$ determines all random draws; $d$ and
  $m$ set the input width and number of target components; $\sigma$ sets the
  observation-noise standard deviation.
  \item[\textbf{Function}] $k_{\min}$, $k_{\max}$, and $p_{\neg}$ determine
  the clause degree bounds and the probability that each selected literal is
  negated.
  \item[\textbf{Distribution structure}] Parent limit $q_{\max}$ and CPT
  range $[p_{\min},p_{\max}]$ bound each input and gate DAG's parents and
  conditional probabilities.
  \item[\textbf{Input dynamics}] $p_x$ controls new input-distribution
  sampling; $p_{\mathrm{sample,min}}^x$ and
  $p_{\mathrm{sample,max}}^x$ bound the shared input-node sampling
  probability.
  \item[\textbf{Gate dynamics}] $p_g$ controls new gate-distribution sampling;
  $p_{\mathrm{sample,min}}^g$ and $p_{\mathrm{sample,max}}^g$ bound the
  shared gate-node sampling probability.
  \item[\textbf{Target scale}] $w_{\min}$, $w_{\max}$, $b_{\min}$,
  $b_{\max}$ bound weights and intercepts; $p_b$ controls new intercept
  sampling.
\end{description}

\noindent\textbf{Initialize:} sample fixed functions $f_i$ and weights $w_i$,
initial intercept $b_0$, input and gate distributions, states $X_0,G_0$, and
their shared node-sampling probabilities $p_{\mathrm{sample}}^x$ and
$p_{\mathrm{sample}}^g$.

\noindent\textbf{For} $t=0,1,\ldots$:
\begin{enumerate}
  \item Set $o_t\gets X_t$ and
  \[
    \tilde y_t\gets b_t+\sum_{i=1}^m w_iG_{t,i}f_i(X_t)+\epsilon_t,
    \qquad \epsilon_t\sim\mathcal N(0,\sigma^2);
  \]
  emit $(o_t,\tilde y_t)$.
  \item Draw $S_t^x$ and $S_t^g$.  If either is one, resample the
  corresponding distribution; retain $X_t$ or $G_t$.
  \item For every input node $j$, draw
  $S_{t,j}^x\sim\operatorname{Bernoulli}(p_{\mathrm{sample}}^x)$; in
  topological order, ancestrally sample nodes with indicator one to produce
  $X_{t+1}$.  Repeat independently
  for every gate node using
  $S_{t,j}^g\sim\operatorname{Bernoulli}(p_{\mathrm{sample}}^g)$ to produce
  $G_{t+1}$.
  \item Draw $S_t^b\sim
  \operatorname{Bernoulli}(p_b)$.  If it is one,
  sample $b_{t+1}\sim\operatorname{Uniform}(b_{\min},b_{\max})$; otherwise
  set $b_{t+1}\gets b_t$.
\end{enumerate}

\section{Drift mechanisms}
Input and gate distribution sampling change their distributions abruptly.
Neither sampling operation directly changes its family's current binary state,
so latent activity continues from its pre-sampling value and is subsequently
updated only by the usual partial ancestral sampling under the new
distribution.  Partial ancestral sampling provides persistent,
within-family correlated states but does not necessarily follow either
configured DAG distribution.  Functions and weights do not drift.
Activity-gate sampling and independent intercept sampling are the sources of output
shifts.
Gates remain unobserved activity variation.  Oracle metadata is optional and
is never supplied to a model.

\section{Multiclass classification task}
The classification variant reuses the observed binary input stream
$x_t\in\{0,1\}^d$ and the same conjunction-of-literals function family
$f_j:\{0,1\}^d\to\{0,1\}$ sampled once at initialization.  For each class
$c\in\{0,\ldots,K-1\}$, the environment maintains hidden activity gates
$g_{t,c}\in\{0,1\}^m$, intercept $b_{t,c}$, and weights $w_{c,j}$ sampled
independently at initialization and fixed thereafter.  The noiseless class
score is
\[
s_{t,c}=b_{t,c}+\sum_{j=1}^{m}w_{c,j}\,g_{t,c,j}\,f_j(x_t).
\]
Independent Gaussian noise perturbs each score,
$\tilde s_{t,c}=s_{t,c}+\epsilon_{t,c}$ with
$\epsilon_{t,c}\sim\mathcal N(0,\sigma^2)$, and the emitted label is
\[
y_t=\operatorname*{argmax}_{c}\tilde s_{t,c},
\]
with ties broken toward the smallest class index.  Input drift matches the
regression task.  Gate distribution resampling, partial gate resampling, and
intercept resampling are independent per class using the same probabilities
$p_g$, $p_{\mathrm{sample,min}}^g$, $p_{\mathrm{sample,max}}^g$, $p_b$,
$b_{\min}$, and $b_{\max}$ as in regression.

\paragraph{Stream generation.}
Initialize shared functions, input distribution, and per-class gate
distributions, gate states, intercepts, and weight matrices.  At each time
step emit $(x_t,y_t)$ from the current states and noisy scores, then apply the
shared input update and independent per-class gate and intercept updates
described above.

\section{API and integration boundaries}
Swappable generators, models, and metrics are required for fair prequential
benchmarks; coupling them would bake evaluation assumptions into model code.
The \texttt{binaml.environments} package owns generation, configuration,
serialization, and state restoration through separate regression and
classification modules.  \texttt{binaml.models} owns online models and has no
environment dependency.  A model implements
\texttt{predict(features)} and \texttt{update(target)}.  The
prediction call retains the features and any derived state;
\texttt{update} applies the revealed target to that stored prediction.
The \texttt{binaml.evaluation}
package combines the two through predict-then-update prequential evaluation.
Named benchmark scenarios only compose these
independent components.

\section{Reference online baselines}
The library provides generic online baselines alongside the binary-feature
models.  For regression, the linear SGD baseline predicts with an affine
function of the observed coordinates and the MLP baseline is a JAX
one-hidden-layer ReLU network.  For classification, the
linear SGD baseline uses softmax cross-entropy and
the MLP baseline uses the same JAX replay scaffold with a softmax head.
JAX baselines are optional and are installed through the \texttt{benchmarks}
dependency extra so ordinary library users do not need JAX.

All provided models use the same replay convention.  After updating with a target,
they retain the most recent $B$ labeled samples and make $S$ optimization
updates on that window.  Linear and MLP baselines both apply an averaged
full-batch gradient update at each step.  Reported comparisons must state each
baseline's architecture, learning rate, regularization, replay size $B$, step
count $S$, optimizer settings, and random seed.

\section{Determinism and reproducibility}
Every trajectory is determined by a validated JSON configuration and a single
root seed $\xi$.  The implementation splits that seed into stable PCG64DXSM
substreams for initialization, sampling, drift, and noise.  Snapshots include
the full latent state and bit-generator states.  Scenarios, generator versions,
and configuration fingerprints are recorded with run artifacts.

\section{Prequential evaluation protocol}
For each emitted sample, an evaluator first records the model prediction, then
exposes the target through \texttt{update}.  Regression benchmarks compute
mean squared error from the recorded pre-update predictions; classification
benchmarks compute accuracy from the recorded pre-update class predictions.
Benchmarks publish their horizon, scenario configuration, seeds, generator
version, and separate total, prediction, and update timings; aggregate
metrics alone are insufficient for replication.

\paragraph{Warm-up.}
A scenario may specify a non-negative \texttt{warmup\_samples} count smaller
than its horizon.  Models predict and update these initial samples normally,
but their errors and step timings are excluded from reported regression MSE
and RMSE or classification accuracy, and from published timings.
Time-series plots retain the full trajectory and
mark the first evaluated sample with a vertical line so that the learning period
remains visible.

\section{Limitations and scope}
These tasks are designed for controlled adaptation studies.  They do not claim
that synthetic drift matches a particular deployment, nor does this document
make comparative performance claims.  Results may be added only for versioned
scenarios with repeatable baseline runs.
\end{document}
